\section{Introdução}
\label{sec:introducao}

As mudanças climáticas figuram entre os temas mais relevantes da ciência contemporânea, e
o estudo de suas tendências depende da análise de séries históricas extensas de
temperatura. Conjuntos de dados como o \emph{Climate Change: Earth Surface Temperature
Data}, da Berkeley Earth~\citep{databerkeley}, reúnem medições de mais de dois séculos por
cidade e país, alcançando milhões de registros --- um volume que torna inviável o
processamento em uma única máquina por ferramentas tradicionais. Esse cenário motiva o uso
de plataformas de \emph{Big Data} capazes de distribuir o cálculo por um aglomerado de
nós.

Este trabalho utiliza o \textbf{Apache Spark}~\citep{sparkdocs}, motor de processamento
distribuído em memória, para investigar tendências de aquecimento e anomalias térmicas ao
longo dos últimos 250 anos. O Spark foi implantado em um \emph{cluster} no modelo
\emph{master/worker}, totalmente conteinerizado com Docker e \texttt{docker-compose}, o que
garante um ambiente reprodutível e isolado. À base de temperaturas somou-se a base de
emissões de CO\textsubscript{2} da \emph{Our World in Data}~\citep{dataowid}, permitindo
correlacionar aquecimento e emissões por meio de um \emph{join} entre as duas fontes.

O objetivo é demonstrar, na prática, a capacidade do Spark de lidar com transformações em
larga escala e cálculos estatísticos complexos, respondendo a oito perguntas analíticas
que envolvem agregações por década, rankings de instabilidade climática, correlações,
\emph{Window Functions} e previsão com regressão linear via \texttt{Spark MLlib}. Ao longo
do \emph{pipeline}, explora-se também o processo de \textbf{ETL} (limpeza, padronização e
enriquecimento dos dados) e técnicas de otimização como o uso de \texttt{cache} para
reaproveitamento de resultados intermediários.

O restante do relatório está organizado da seguinte forma: a Seção~\ref{sec:arquitetura}
descreve a arquitetura do \emph{cluster}; a Seção~\ref{sec:etl} detalha a metodologia de
ETL; a Seção~\ref{sec:cache} discute a análise de otimização com \texttt{cache}; a
Seção~\ref{sec:respostas} apresenta as respostas às oito perguntas com suas visualizações e
interpretações; e a Seção~\ref{sec:sparkui} reúne as evidências de execução coletadas na
Spark UI.
